% TOP_PROBLEM: {"id": 1462, "title": "Tarski's Exponential Function Problem", "category_id": 4, "category": {"id": 4, "name": "algebra", "display_name": "Algebra", "description": "Group theory, ring theory, field theory, and algebraic structures.", "slug": "algebra", "order_index": 4, "created_at": "2026-07-31T15:26:25.671Z"}, "status": "open"}
% FIRST_POSTED: 2026-09-27
% !TeX program = lualatex
% Compile twice: lualatex 1462.tex
% All references and prose are embedded; no BibTeX database or external inputs.
\documentclass[11pt,a4paper]{article}
\usepackage[margin=24mm,headheight=14pt]{geometry}
\usepackage{fontspec}
\setmainfont{Latin Modern Roman}
\setsansfont{Latin Modern Sans}
\setmonofont{Latin Modern Mono}
\usepackage{amsmath,amssymb,mathtools}
\usepackage{unicode-math}
\setmathfont{Latin Modern Math}
\usepackage{microtype}
\usepackage{enumitem,xurl,needspace}
\usepackage[dvipsnames]{xcolor}
\usepackage{hyperref}
\hypersetup{unicode=true,colorlinks=true,linkcolor=MidnightBlue,urlcolor=MidnightBlue,
 pdftitle={500 Mathematical Problem Statements},pdfauthor={Source-annotated compilation},bookmarksnumbered=true}
\usepackage{bookmark}
\usepackage{tocloft}
\setlength{\cftsecnumwidth}{3em}
\cftsetpnumwidth{2.5em}
\cftsetrmarg{4em}
\usepackage{titlesec}
\titleformat{\section}{\Large\bfseries\raggedright}{\thesection}{0.7em}{}
\usepackage{fancyhdr}
\pagestyle{fancy}\fancyhf{}
\fancyhead[L]{\small\sffamily Mathematical problem statements}
\fancyhead[R]{\small Source edition 22}
\fancyfoot[C]{\thepage}
\setlength{\parindent}{0pt}
\setlength{\parskip}{5pt plus 1pt}
\setlength{\emergencystretch}{3em}
\setcounter{tocdepth}{1}
\setcounter{secnumdepth}{1}
\setlist[itemize]{leftmargin=3em,itemsep=5pt,topsep=4pt,parsep=0pt}
\allowdisplaybreaks
\newcommand{\N}{\mathbb{N}}
\newcommand{\Z}{\mathbb{Z}}
\newcommand{\Q}{\mathbb{Q}}
\newcommand{\R}{\mathbb{R}}
\newcommand{\C}{\mathbb{C}}
\newcommand{\F}{\mathbb{F}}
\newcommand{\A}{\mathbb{A}}
\newcommand{\PP}{\mathbb P}
\newcommand{\E}{\mathbb{E}}
\newcommand{\eps}{\varepsilon}
\newcommand{\dd}{\,\mathrm d}
\newcommand{\sref}[2]{\hyperlink{source:#1:#2}{[S#2]}}
\newcommand{\eref}[2]{\hyperlink{extra:#1:#2}{[E#2]}}
% Turn the next switch off for a shorter reading copy. The quotations preserve
% every exactTarget field, including qualifications omitted from a short summary.
\newif\ifshowcatalogue
\showcataloguetrue
\newcommand{\cataloguescope}[1]{\ifshowcatalogue
 \paragraph{Exact scope in the supplied catalogue.}
 \begin{quote}\small #1\end{quote}\fi}
\hypersetup{pdftitle={Problem 1462 - Tarski's exponential problem},pdfauthor={Open Problems Math research notebook}}
\fancyhead[L]{\small\sffamily Open Problems Math \textbar\ Research notebook}
\fancyhead[R]{\small\sffamily 124 \textbar\ 27 September 2026}
\numberwithin{equation}{section}
\begin{document}
\setcounter{section}{0}
% ================================================================
% problemId: 1462
% Canonical problem ID: 1462
% exactTarget SHA-256: 62800700ac4f4025ab695be687817d0d33fbdcef617bff09c427d5f30465a19a
\clearpage
\section*{\texorpdfstring{Tarski's Exponential Function Problem}{Tarski's Exponential Function Problem}}\label{1462}
\subsection{Definitions and mathematical statement}
A first-order sentence in the real exponential field is built from real variables, the constants $0,1$, addition, multiplication, order, equality, and the usual real function $\exp$, using Boolean operations and quantifiers over $\R$. Let
\[
 \operatorname{Th}(\R_{\exp})=\{\ulcorner\varphi\urcorner:(\R;0,1,+,\cdot,\exp,<)\models\varphi\}.
\]
The question is whether this set of finite encodings is decidable by an algorithm halting on every input. It concerns arbitrary quantifier alternation, not only an existential fragment, and imposes no complexity bound.
\par\smallskip\textit{Formulation and concept references:} \sref{1462}{1}.
\cataloguescope{Determine whether the first-order theory of the real exponential field (R; 0, 1, +, *, exp, <, =) is decidable.}
\subsection{Short English statement}
Can a terminating algorithm decide every first-order statement about real arithmetic with the exponential function?
% BEGIN RESEARCH ADDITION 124
\subsection{Research attempt}

\paragraph{Current frontier.}
The supplied STACS 2025 paper distinguishes Tarski's full exponential
field problem from its own existential extension by integer powers.
Its introductory root-barrier conjecture gives a useful formulation of
the effective zero-testing obstacle. \sref{1462}{1}, Section 1,
Conjecture 2.
Berarducci--Gallinaro's June 2026 revision gives a new axiomatization and
decidability argument \emph{assuming the real form of Schanuel's
conjecture}; an unconditional restricted model-completeness theorem in
that paper is not an unconditional solution of the full decision
problem. \eref{1462}{1}, introduction and Theorems 12.4 and 14.4.
The following attempt separates effective transformations actually
available from transcendence hypotheses not proved here.

\paragraph{Attack routes.}
One may algebraize exponential terms by changing variables, use certified
numerical root isolation, or enumerate consequences of an effective
complete axiomatization.  The first route must cope with variables
appearing both outside and inside exponentials.  The second needs an
effective gap between zero and a small nonzero value, including singular
roots.  The third needs completeness, not just a recursively enumerable
sound list of axioms.  We develop an exact algebraizable fragment and
then examine the missing zero test.

\paragraph{Attempt I: an all-quantifier decidable fragment.}
Consider sentences whose atomic expressions are rational polynomials in
terms
\[
 \exp(r_1x_1+\cdots+r_sx_s),\qquad r_i\in\Q,
\]
with arbitrary Boolean combinations and quantifier alternation.  Require
that the variables occur \emph{only} in these homogeneous rational
linear exponents.  There are no bare $x_i$, no nested exponentials, and
no nonzero constant terms inside the exponent.  These restrictions
define a fragment of the target language, not a replacement for it.

Take a positive integer $D$ divisible by all coefficient denominators
and substitute $y_i=\exp(x_i/D)$.  This is a bijection from $\R$ to
$\R_{>0}$ in each coordinate.  If $r_i=a_i/D$, then
\[
 \exp\!\left(\sum_i r_i x_i\right)=\prod_i y_i^{a_i}.
\]
Every atom becomes a comparison of rational Laurent polynomials.
Multiplying both sides of an atomic equality or inequality by a
sufficiently large positive monomial clears all negative exponents
without changing its truth value.  Replace quantifiers by
\[
 \exists x_i\,\varphi\ \rightsquigarrow\
 \exists y_i\,(y_i>0\wedge\varphi^*),\qquad
 \forall x_i\,\varphi\ \rightsquigarrow\
 \forall y_i\,(y_i>0\Rightarrow\varphi^*).
\]
By structural induction on the formula, this transformation preserves
truth, including every quantifier alternation.  It terminates after
finite rational arithmetic.  The resulting sentence lies in the theory
of real closed fields, whose decision procedure is the classical
Tarski result recalled in \eref{1462}{1}, Section 1.

For example, the sentence
$\forall x\,\exists z\,[e^{x/2}e^z=1]$ becomes
$\forall y>0\,\exists w>0\,[yw^2=1]$ with $D=2$ and is true.
The equally elementary-looking equation $e^x=x+2$ lies outside the
fragment: after the substitution, the right side contains $D\log y+2$.
Introducing a symbol for $\log$ merely restores the original
transcendental relation.  Adding arbitrary exponential constants also
requires deciding their exact algebraic relations, which the
transformation did not do.

\paragraph{Attempt II: what numerical certification can decide.}
Let $f_1,\ldots,f_r$ be specified exponential-polynomial expressions on
a compact rational box $B$.  Bounds on their derivatives can be
computed recursively from rational arithmetic and upper bounds on the
exponential function on bounded intervals.  Consequently
$F(x)=\min_i f_i(x)$ has a computable Lipschitz constant $L$.
For a supplied rational $\delta>0$, the following \emph{promise problem}
is decidable:
\[
 \max_{x\in B}F(x)\geq\delta
 \quad\hbox{versus}\quad
 \max_{x\in B}F(x)\leq-\delta.
\]
Choose a finite rational grid with covering radius at most
$\delta/(4\max(1,L))$ and evaluate each $F$ to error at most
$\delta/4$.  The grid maximum differs from the true maximum by at
most $\delta/2$.  Its sign distinguishes the two cases.  Every
operation terminates; the number of grid points can be large but the
original task imposes no complexity bound.

This procedure is not a solution of the unrestricted existential
problem.  If the maximum is zero, no fixed positive promise gap applies.
Nor does it decide unbounded searches by checking successively larger
boxes: an empty set provides no automatic stopping certificate.
For $f(x)=e^x-x-1$, the exact zero at $x=0$ is singular since $f'(0)=0$.
Interval errors getting smaller around that value do not alone certify
whether a perturbed expression is exactly zero.

The root-barrier assertion in the formulation source is more targeted.
Given integer exponential polynomials $f_1,\ldots,f_n,g$, it asks for
an effective positive integer $t$ such that at every nonsingular real
zero $\alpha$ of the $f_i$ one has either $g(\alpha)=0$ or
$|g(\alpha)|>1/t$. \sref{1462}{1}, Conjecture 2.
If this $t$ and certified arbitrarily accurate approximations to an
isolated such $\alpha$ were available, evaluating $g(\alpha)$ within
$1/(4t)$ would distinguish the two alternatives by the threshold
$1/(2t)$.  This implication is elementary and exact.  The unproved
part is an algorithm producing the uniform $t$ with the source's
quantifiers, not the final comparison.  The presence of only finitely
many values in some isolated calculation would not supply such an
effective bound uniformly across inputs.

\paragraph{Attempt III: testing the proof-enumeration route.}
Suppose $T$ were a recursively enumerable, sound, \emph{complete}
axiomatization of $\operatorname{Th}(\R_{\exp})$.  Dovetail proof
searches for $\varphi$ and $\neg\varphi$.  Soundness prevents both from
appearing, and completeness guarantees that one appears, yielding a
decision procedure.  If completeness is omitted, neither search need
halt.  Thus writing down plausible differential axioms for an ordered
exponential field does not establish decidability.

The recent paper proves the required completeness for its indicated
axioms only under $\mathrm{SC}_{\R}$, the real Schanuel assertion about
transcendence degree of linearly independent real tuples and their
exponentials.  That hypothesis has not been proved in this notebook.
The paper's unconditional model completeness in a restricted setting
cannot be substituted for the full completeness statement.
\eref{1462}{1}, Theorems 12.4, 14.2 and 14.4.

\paragraph{Stress test.}
The algebraic change of variables is globally bijective on each
quantified coordinate; it is not just a local analytic substitution.
Cleared denominators are positive, so inequality directions are
preserved.  The fragment excludes exactly the mixed ordinary/exponential
occurrences that would reintroduce logarithms.  The finite-grid algorithm
states its promise and compact domain explicitly, rather than mistaking
a semidecision procedure for a halting algorithm on every sentence.
The root-barrier discussion includes nonsingularity and an effective
bound; qualitative isolation alone is not enough.

The companion symbolic checks verify Laurent identities and positive
monomial clearing on examples.  They are not an implementation of a
full real-exponential decision procedure or a test of Schanuel's
conjecture.  No finite precision convention is allowed to identify a
small value with zero in the mathematical argument.

\paragraph{Outcome.}
\textbf{Outcome: Exploratory approach with unresolved gap.}

An explicit all-quantifier fragment has been reduced effectively to
real closed fields, and a robust compact-domain promise problem has a
proved grid algorithm.  Neither includes arbitrary mixed exponential
sentences.  The precise missing zero-barrier/completeness inputs have
been isolated, with their dependence on an unresolved transcendence
hypothesis retained.  The original decidability question is not resolved,
and no novelty claim is made for these fragment arguments.
% END RESEARCH ADDITION 124
\subsection{Sources}
\begin{itemize}\raggedright
\item[\textnormal{[S1]}]\hypertarget{source:1462:1}{} Jorge Gallego-Hernandez, Alessio Mansutti, "On the Existential Theory of the Reals Enriched with Integer Powers of a Computable Number," 42nd International Symposium on Theoretical Aspects of Computer Science (STACS 2025), LIPIcs vol. 327, Article 37, pp. 37:1-37:18.
\url{https://drops.dagstuhl.de/storage/00lipics/lipics-vol327-stacs2025/LIPIcs.STACS.2025.37/LIPIcs.STACS.2025.37.pdf}.
{\small\textit{Catalogue formulation source.}}
% BEGIN RESEARCH REFERENCES 124
\item[\textnormal{[E1]}]\hypertarget{extra:1462:1}{} Alessandro Berarducci and Francesco Gallinaro.
\emph{On the elementary theory of the real exponential field}, arXiv:2603.08365v2, 23 June 2026. Theorems 12.4, 14.2 and 14.4. Restricted model completeness is unconditional; the full completeness and decidability conclusions retain the real Schanuel hypothesis.
\url{https://arxiv.org/html/2603.08365v2}.
{\small\textit{Consulted 27 September 2026 for the indicated result or scope.}}
% END RESEARCH REFERENCES 124
\end{itemize}


\end{document}
